% II. Trabajos relacionados
\subsection{Conjuntos de datos de RM lumbar}
El \emph{Lumbar Spine MRI Dataset} reúne estudios de 515 pacientes con dolor lumbar sintomático \cite{sudirman2019}. Sobre ese mismo corpus, Natalia \emph{et al.} seleccionaron los 515 cortes T2 medio-sagitales y predijeron el grado de Pfirrmann con una exactitud media del 88,1\,\% \cite{natalia2024}, y Masood \emph{et al.} publicaron 1\,545 grados de Pfirrmann validados por un radiólogo para L3-L4, L4-L5 y L5-S1 \cite{masood2026}. Fuera de este corpus, SPIDER ofrece 447 series sagitales T1 y T2 de 218 pacientes de cuatro hospitales, anotadas de forma iterativa \cite{vandergraaf2024}, y LumbarDISC reúne 2\,697 pacientes de ocho instituciones de seis países, anotados por radiólogos expertos \cite{richards2026}. Los conjuntos multicéntricos y anotados por expertos son la referencia para validar modelos, pero el de Sudirman sigue siendo uno de los pocos que permiten trabajar con discos individuales en un plazo académico.

\subsection{Clasificación de discos con aprendizaje profundo}
La familia YOLO es la más usada en los trabajos revisados. Liawrungrueang \emph{et al.} detectaron y graduaron la degeneración con YOLOv5 sobre 1\,000 imágenes T2 sagitales de 515 pacientes, con exactitud superior al 95\,\% \cite{liawrungrueang2023}; Prisilla \emph{et al.} compararon YOLOv5x, YOLOv6 y YOLOv7 para detectar hernias con etiquetas de radiólogo \cite{prisilla2023}; Wang \emph{et al.} graduaron degeneración y hernia con un YOLO modificado en T2 sagital y axial, con 420 pacientes internos y 52 externos \cite{wang2025}. Sayed \emph{et al.} usaron YOLOv8 para detectar los discos y recortar un ROI por disco en cortes sagitales (mAP50 del 99,4\,\% para el disco) \cite{sayed2025scirep}, y un flujo en cascada de detector y CNN para identificar discos y detectar hernias en cortes axiales y sagitales \cite{sayed2025imu}. Con redes residuales, Tan \emph{et al.} combinaron ResNet-50 y YOLOv2 para clasificar discos sanos y herniados, aunque con solo 32 pacientes (160 discos) \cite{tan2024}, y Al-kubaisi y Khamiss aplicaron transferencia de aprendizaje a las mismas cuatro clases de este proyecto en 181 pacientes, sobre todo en cortes axiales \cite{alkubaisi2022}. En tomografía computarizada, abultamiento y protrusión fueron el par de clases más confundido \cite{gao2025}.

\subsection{Características profundas con clasificadores clásicos}
Sobre el mismo corpus de Sudirman, Sarı y Atasoy extrajeron características con ResNet-50, las seleccionaron con optimización por lobo gris (GWO) y las clasificaron con SVM, con un 99,42\,\% en una tarea binaria con cortes T1 y validación cruzada de cinco particiones, sin declarar si la división fue por paciente \cite{sari2025}. Khoiwal y McMillan mostraron que clasificadores lineales sobre \emph{embeddings} congelados igualan a modelos entrenados desde cero en retina, mamografía, dermatoscopia y tórax \cite{khoiwal2024}. En datos tabulares de la columna (ángulos óseos), el mejor clasificador fue un SMOTE-RFE-XGBoost \cite{zhang2023}.

\subsection{Vacío de investigación}
Los resultados más altos de la literatura (97--99\,\%) corresponden a tareas binarias o a conjuntos de prueba pequeños, y las cifras internas bajan en validación externa: en un estudio, el F1 sagital pasó de 0,971 a 0,768 y 0,787 en dos centros externos \cite{yi2023}. Pocos trabajos sobre el corpus de Sudirman declaran una división por paciente, reportan métricas por clase con el desbalance real o analizan sesgos como el nivel vertebral. Este trabajo aborda ese vacío: compara cuatro familias de modelos sobre la misma división por paciente, evalúa sobre imágenes originales y reporta métricas por clase y por nivel.
